\scanpage{49}
séparé $E\otimes F$, \deleted{\unsure{dont}} \inserted{\unsure{où}} le produit scalaire est défini par (3).
La norme d'un élément \deleted{\illegible{2}} $u\in E\overset{(2)}\otimes F$ est notée $\|u\|_2$.
$E\overset{(2)}\otimes F$ est donc un espace de Hilbert.

\uline{Prop. 1} Soient $E$ et $F$ deux espaces de Hilbert. \deleted{\unsure{$x\in E$, $y\in F$,}}
\deleted{$\|x\otimes y\|_2=\|x\|\|y\|\unsure{{}^2}$}
Soit $(e_i)_{i\in I}$ une base orthonormale dans $E$, $(f_j)_{j\in J}$ une base orthonormale dans $F$, alors $(e_i\otimes f_j)_{(ij)\in I\times J}$ est une base orthonormale dans $E\overset{(2)}\otimes F$.

En effet, d'après (4) le système $(e_i\otimes f_j)$ est orthonormal, \deleted{de plus les \unsure{combinaisons finies sont denses dans} $E\overset{(2)}\otimes F$.}
\deleted{En effet, \illegible{2}} \unsure{montrons} que l'espace vectoriel qu'il engendre est dense dans $E\overset{(2)}\otimes F$.
\deleted{Pour ceci}, si $E_0$ (resp. $F_0$) \deleted{est} \inserted{désigne} l'espace engendré par les $e_i$ (resp. les $f_j$) \deleted{\unsure{et les combinaisons finies}} \deleted{\unsure{par suite}}
$E_0$ est dense dans $E$, $F_0$ dense dans $F$, donc $E_0\otimes F_0$ dense dans $E\otimes F$, donc aussi dans $E\overset{(2)}\otimes F$.

Notons aussi la formule
\[
(7)\qquad\boxed{\|x\otimes y\|_2=\|x\|\|y\|}
\]
\unsure{résulte} en effet \deleted{de} \deleted{qui résulte aussitôt de}
\[
(x\otimes y,x\otimes y)=(x,x)(y,y).
\]
\deleted{\unsure{le cas de forme} \illegible{3}}
Notons aussi que l'isomorphisme \inserted{canonique} $\overline{E\otimes F}=\overline E\otimes\overline F$ (formule (\unsure{2|1})) conserve évidemment les structures préhilbertiennes, d'où un isomorphisme canonique
\[
(8)\qquad\overline{E\overset{(2)}\otimes F}=\overline E\overset{(2)}\otimes\overline F.
\]
\mspage{50}
